% SOURCE_STATEMENT: "Does the Asaeda–Haagerup subfactor arise from a conformal field theory?"
% Preparation and final self-audit: CONTRIBUTING.md prompts C and D; source/audit artifacts retained in the local physics_statements_2026_10_11 batch.
% !TeX program = lualatex
\documentclass[11pt]{article}
\usepackage[margin=25mm]{geometry}
\usepackage{fontspec}
\setmainfont{Latin Modern Roman}
\usepackage{amsmath,amssymb,mathtools}
\usepackage{unicode-math}
\setmathfont{Latin Modern Math}
\usepackage{xurl,hyperref}
\hypersetup{colorlinks=true,urlcolor=blue}
\setcounter{secnumdepth}{0}
\setlength{\parindent}{0pt}
\setlength{\parskip}{5pt}
\setlength{\emergencystretch}{3em}
\newcommand{\sref}[2]{\hyperlink{source:#1:#2}{[S#2]}}
\newcommand{\cataloguescope}[1]{\paragraph{Recorded target.}#1}
\newcommand{\R}{\mathbb R}
\newcommand{\C}{\mathbb C}
\newcommand{\N}{\mathbb N}
\newcommand{\Z}{\mathbb Z}
\newcommand{\Q}{\mathbb Q}
\newcommand{\F}{\mathbb F}
\newcommand{\eps}{\varepsilon}
\begin{document}
% UnsolvedMath ID 30002840; source catalogue code OWR-13673-010
\setcounter{section}{30002839}
\section{Conformal Field Theory Realization of the Asaeda–Haagerup Subfactor}\label{30002840}
{\small UnsolvedMath ID 30002840 · Mathematical Physics}
\subsection{Definitions and mathematical statement}

\paragraph{Shared notation.}$\mathbb N=\{1,2,\ldots\}$; the standard number fields have their usual meanings. All additional conventions are stated locally.


\paragraph{Definitions and conventions.}
A subfactor is an inclusion $N\subset M$ of von Neumann algebras with trivial centers. Its index and standard invariant encode the size and higher relative-commutant structure of the inclusion. In an operator-algebraic conformal field theory, interval or local observable algebras and sectors can produce subfactor inclusions. To say a subfactor arises in conformal field theory requires identifying the inclusion or invariant that the construction realizes. The Asaeda--Haagerup subfactor is the exceptional finite-index subfactor named in the source; its standard invariant is the specific datum whose conformal realization is sought.

\paragraph{Target retained from the canonical source catalogue.}
\paragraph{Catalogue wording retained for traceability.}
Does the Asaeda–Haagerup subfactor arise from a conformal field theory?

\paragraph{Source scope and qualification.}
This is a named special case of the general realization question, IDs 30002833 and 30004278. An unrelated subfactor with similar index does not realize the target standard invariant.
\subsection{Short English statement}
Can conformal field theory realize the exceptional Asaeda--Haagerup subfactor?
\subsection{Sources}
\begin{itemize}
\item[\hypertarget{source:30002840:1}{S1}] ULAM.AI and the UnsolvedMath Contributors, UnsolvedMath, record 30002840 (OWR-13673-010). Catalogue curation under CC BY 4.0; definitions and English exposition added locally.
\url{https://huggingface.co/datasets/ulamai/UnsolvedMath}.
\item[\hypertarget{source:30002840:2}{S2}] Subfactors and Conformal Field Theory (2015), Oberwolfach Reports, DOI 10.4171/owr/2015/16, pp. 849--926. Primary report accessed 2026-10-11; relevant PDF pages 35, 36.
\url{https://doi.org/10.4171/owr/2015/16}.
\item[\hypertarget{source:30002840:3}{S3}] Grossman--Izumi--Snyder, Asaeda--Haagerup categories (catalogue-listed context; not independently consulted in this preparation).
\url{https://arxiv.org/abs/1501.07324}.
\end{itemize}
\label{end:30002840}
\end{document}
